% ====================================================================
%  徐森林、金亚东、薛春华《数学分析》（第三册）
%  复习题 16
%  源：原书印刷页 305–308（PDF 物理页 311–314）
%  约定：只用 exblock，不设 section 命令（故不入目录）
% ====================================================================

% ----------------------------------------------------------------------
% 复习题 16
% ----------------------------------------------------------------------
\begin{exblock}{复习题 16}
\textbf{1.} 设 $g(t)$ 为区间 $[0,h]$ ($h>0$) 上的单调增函数, 利用
\[
  \int_{0}^{h}g(t)\frac{\sin\lambda t}{t}\mathrm{d}t=g(0+0)\int_{0}^{h}\frac{\sin\lambda t}{t}\mathrm{d}t+\int_{0}^{h}\left[g(t)-g(0+0)\right]\frac{\sin\lambda t}{t}\mathrm{d}t
\]
证明:
\[
  \lim_{\lambda\to+\infty}\int_{0}^{h}g(t)\frac{\sin\lambda t}{t}\mathrm{d}t=\frac{\pi}{2}g(0+0).
\]

\textbf{2.} 设 $0<x<2\pi$, 证明:
\begin{enumerate}
\item $\displaystyle\sum_{k=1}^{n}\frac{\sin kx}{k}=-\frac{x}{2}+\int_{0}^{x}\frac{\sin\left(n+\frac{1}{2}\right)t}{2\sin\frac{t}{2}}\mathrm{d}t$;
\item $\displaystyle\sum_{k=1}^{\infty}\frac{\sin kx}{k}=\frac{\pi-x}{2}$.
\end{enumerate}

\textbf{3.} 利用 $\cos ax$ 在 $[-\pi,\pi]$ 上的 Fourier 展开式证明:
\begin{enumerate}
\item $\displaystyle\cot x=\frac{1}{x}+\sum_{n=1}^{\infty}\frac{2x}{x^{2}-n^{2}\pi^{2}},\quad x\neq k\pi,\ k=0,\pm1,\pm2,\cdots$;
\item $\displaystyle\frac{1}{\sin x}=\frac{1}{x}+\sum_{n=1}^{\infty}(-1)^{n}\frac{2x}{x^{2}-n^{2}\pi^{2}},\quad x\neq k\pi,\ k=0,\pm1,\pm2,\cdots$.
\end{enumerate}

\textbf{4.} 试用 Fej\'er 定理~\ref{thm:16.4.2} 证明: 如果周期为 $2\pi$ 的连续函数 $f$ 与三角函数系
\[
  \{1,\cos x,\sin x,\cdots,\cos nx,\sin nx,\cdots\}
\]
中每个函数都正交, 则必有 $f=0$.

\textbf{5.} 定义在区间 $[0,1]$ 上的函数系
\[
  \varphi_{n}(t)=\operatorname{sgn}(\sin 2^{n}\pi t),\quad n=1,2,\cdots
\]
称为 \textbf{Rademacher 函数系}. 证明: Rademacher 函数系是 $[0,1]$ 上的规范正交系.

\textbf{6.} 设 $\displaystyle\sum_{n=0}^{\infty}a_{n}$, $\displaystyle\sum_{n=0}^{\infty}b_{n}$ 为两个收敛级数, 其和分别为 $A$, $B$. 证明: 它们的 Cauchy 乘积 $\displaystyle\sum_{n=0}^{\infty}c_{n}$ 必能 Abel 求和, 且
\[
  \sum_{n=0}^{\infty}c_{n}=AB(\mathrm{A}).
\]

\textbf{7.} 试由 Weierstrass 的三角多项式逼近定理导出代数多项式逼近定理.

\textbf{8.} 设 $f(x)$ 为 $[0,\pi]$ 上的二阶连续可导函数, 且 $f(0)=f(\pi)=0$, 令 $f(x)$ 的 Fourier 展开式为
\[
  f(x)=\sum_{k=1}^{\infty}a_{k}\sin kx,
\]
其部分和为 $S_{n}(x)=\displaystyle\sum_{k=1}^{n}a_{k}\sin kx$, $n=1,2,\cdots$, 证明:
\begin{enumerate}
\item $\displaystyle a_{k}=-\frac{2}{\pi k^{2}}\int_{0}^{\pi}f''(x)\sin kx\,\mathrm{d}x,\quad k=1,2,\cdots$;
\item $\displaystyle\int_{0}^{\pi}[f(x)-S_{n}(x)]^{2}\mathrm{d}x\leqslant\frac{1}{3n^{3}}\int_{0}^{\pi}[f''(x)]^{2}\mathrm{d}x,\quad n=1,2,\cdots$.
\end{enumerate}

\textbf{9.} 设 $f(\theta)$ 为 $\mathbb{R}$ 上的周期为 $2\pi$ 的连续函数, 且
\[
  f(\theta)\sim\frac{a_{0}}{2}+\sum_{n=1}^{\infty}\left(a_{n}\cos n\theta+b_{n}\sin n\theta\right).
\]
证明: (1) $\displaystyle u_{n}=\frac{a_{0}}{2}+\sum_{k=1}^{n}r^{k}\left(a_{k}\cos k\theta+b_{k}\sin k\theta\right)$ 在单位圆盘 $D=\{z\in\mathbb{C}\ (\text{复数域})\mid |z|<1\}$ 内的紧子集上一致收敛于一个调和函数 $u(x,y)$ (即 $\Delta u=\dfrac{\partial^{2}u}{\partial x^{2}}+\dfrac{\partial^{2}u}{\partial y^{2}}=0$), 其中 $z=re^{\mathrm{i}\theta}=x+\mathrm{i}y$, $\mathrm{i}=\sqrt{-1}$;
\begin{enumerate}
\item[(2)] $\displaystyle\iint_{D}\left[\left(\frac{\partial u}{\partial x}\right)+\left(\frac{\partial u}{\partial y}\right)^{2}\right]\mathrm{d}x\,\mathrm{d}y=\pi\sum_{n=1}^{\infty}n\left(a_{n}^{2}+b_{n}^{2}\right)$.
\end{enumerate}

\textbf{10.} 证明: Legendre 多项式
\[
  X_{n}(x)=\frac{1}{2^{n}n!}\frac{\mathrm{d}^{n}(x^{2}-1)^{n}}{\mathrm{d}x^{n}}=\frac{1}{(2n)!!}\frac{\mathrm{d}^{n}(x^{2}-1)^{n}}{\mathrm{d}x^{n}}
\]
在 $[-1,1]$ 上为正交系, 并且 $\displaystyle\int_{-1}^{1}X_{n}^{2}(x)\mathrm{d}x=\frac{2}{2n+1}$.

\textbf{11.} 设函数列 $\{y_{n}(x)\}$ 满足方程
\[
  \frac{\mathrm{d}}{\mathrm{d}x}\left[p(x)\frac{\mathrm{d}y_{n}}{\mathrm{d}x}\right]+\lambda_{n}y_{n}=0,\quad\forall x\in[a,b],\ n=1,2,\cdots,
\]
其中 $n\neq m$ 时, $\lambda_{n}\neq\lambda_{m}$, 且有边界条件 $y_{n}(a)=y_{n}(b)=0$. 证明: $\{y_{n}(x)\}$ 为 $[a,b]$ 上的正交系.

\textbf{12.} 设 $0<\lambda_{1}<\lambda_{2}<\cdots<\lambda_{n}<\cdots$ 满足
\[
  \sigma\sin\sqrt{\lambda_{n}}\,l+\sqrt{\lambda_{n}}\cos\sqrt{\lambda_{n}}\,l=0,\quad\sigma>0.
\]
证明: $\{y_{n}(x)\}=\{\sin\sqrt{\lambda_{n}}\,x\}$ 在 $[0,l]$ 上为正交系, 并求 $\displaystyle\int_{0}^{l}y_{n}^{2}(x)\mathrm{d}x$ (用 $\sigma,\lambda_{n},l$ 表示).

\textbf{13.} 设 $f(x)$ 是以 $2\pi$ 为周期的函数, 满足 $\alpha$ 阶 H\"older 条件(也称 $\alpha$ 阶 Lipschitz 条件):
\[
  |f(x)-f(y)|\leqslant L|x-y|^{\alpha},\quad 0<\alpha\leqslant1.
\]
证明: $\displaystyle a_{n}=O\left(\frac{1}{n^{\alpha}}\right)$, $\displaystyle b_{n}=O\left(\frac{1}{n^{\alpha}}\right)$, $n\to+\infty$.

\textbf{14.} 设 $f(x)$ 有界, 周期为 $2\pi$, 并在 $(-\pi,\pi)$ 上逐段单调. 证明:
\[
  a_{n}=O\left(\frac{1}{n}\right),\quad b_{n}=O\left(\frac{1}{n}\right),\quad n\to+\infty.
\]

\textbf{15.} (1) 证明: 周期为 $2\pi$ 的函数
\[
  f(x)=\frac{1}{4}x(2\pi-x),\quad x\in[0,2\pi]
\]
的 Fourier 级数为
\[
  \frac{1}{4}x(2\pi-x)=\frac{\pi^{2}}{6}-\sum_{n=1}^{\infty}\frac{1}{n^{2}}\cos nx,\quad x\in[0,2\pi].
\]
并由此证得 $\displaystyle\sum_{n=1}^{\infty}\frac{1}{n^{2}}=\frac{\pi^{2}}{6}$.

(2) 应用 (1) 中公式及逐项积分证明:
\begin{align*}
  \frac{1}{6}\pi^{2}x-\frac{1}{4}\pi x^{2}+\frac{x^{3}}{12}&=\sum_{n=1}^{\infty}\frac{1}{n^{3}}\sin nx,\\
  -\frac{1}{36}\pi^{2}x^{3}+\frac{1}{48}\pi x^{4}-\frac{x^{5}}{240}&=\sum_{n=1}^{\infty}\frac{1}{n^{4}}\left(\frac{\sin nx}{n}-x\right),\quad x\in[0,2\pi].
\end{align*}
由此得到 $\displaystyle\sum_{n=1}^{\infty}\frac{1}{n^{4}}=\frac{1}{90}\pi^{4}$.

\textbf{16.} 设 $f(x)$ 是以 $2\pi$ 为周期的连续偶函数, 它的 Fourier 级数为
\[
  f(x)\sim\frac{a_{0}}{2}+\sum_{n=1}^{\infty}a_{n}\cos nx.
\]
证明: (1) 函数 $\displaystyle H(x)=\frac{1}{\pi}\int_{-\pi}^{\pi}f(x+t)f(t)\mathrm{d}t$ 是以 $2\pi$ 为周期的连续偶函数, 它的 Fourier 级数为
\[
  H(x)\sim\frac{a_{0}^{2}}{2}+\sum_{n=1}^{\infty}a_{n}^{2}\cos nx;
\]
(2) $\displaystyle\frac{a_{0}^{2}}{2}+\sum_{n=1}^{\infty}a_{n}^{2}\cos nx$ 一致收敛于 $H(x)$.

\textbf{17.} 设 $f(x)=x$ 在 $(-\pi,\pi)$ 上的 Fourier 级数为
\[
  x=\sum_{n=1}^{\infty}\frac{2(-1)^{n-1}}{n}\sin nx,\quad x\in(-\pi,\pi).
\]
试求函数 $\varphi(x)=x\sin x$ 在 $[-\pi,\pi]$ 上的 Fourier 级数.

\textbf{18.} 设 $f(x)$ 的 Fourier 级数为
\[
  f(x)\sim\frac{a_{0}}{2}+\sum_{n=1}^{\infty}\left(a_{n}\cos nx+b_{n}\sin nx\right).
\]
证明: $f(x)\sin x$ 的 Fourier 级数为
\[
  f(x)\sin x\sim\frac{b_{1}}{2}+\sum_{n=1}^{\infty}\left(\frac{b_{n+1}-b_{n-1}}{2}\cos nx+\frac{a_{n-1}-a_{n+1}}{2}\sin nx\right),
\]
其中 $b_{0}=0$.
\end{exblock}
